% =============================================================================
%  Guion formal de la presentacion: Sueno, memoria y senales fisiologicas (v3)
%  Documento de exposicion. Compilar con pdflatex (dos pasadas por el indice).
%  Nota: las rutas/nombres en \texttt{} van SIN acentos (son archivos reales).
% =============================================================================
\documentclass[11pt,a4paper]{article}

\usepackage[utf8]{inputenc}
\usepackage[T1]{fontenc}
\usepackage[spanish,es-noshorthands]{babel}
\usepackage{lmodern}
\usepackage{microtype}
\usepackage[a4paper,margin=2.3cm]{geometry}
\usepackage{xcolor}
\usepackage{enumitem}
\usepackage{booktabs}
\usepackage{array}
\usepackage{longtable}
\usepackage{colortbl}
\usepackage{titlesec}
\usepackage{fancyhdr}
\usepackage[hidelinks]{hyperref}

% --- Paleta -----------------------------------------------------------------
\definecolor{navy}{HTML}{1F3B63}
\definecolor{steel}{HTML}{3F72AF}
\definecolor{amber}{HTML}{E08A2E}
\definecolor{ink}{HTML}{1E293B}
\definecolor{muted}{HTML}{6B7280}
\definecolor{warn}{HTML}{B4451F}
\definecolor{good}{HTML}{2E7D5B}
\color{ink}

% --- Titulos ----------------------------------------------------------------
\titleformat{\section}{\Large\bfseries\color{navy}}{\thesection.}{0.5em}{}
\titleformat{\subsection}{\large\bfseries\color{steel}}{\thesubsection}{0.5em}{}
\titlespacing*{\section}{0pt}{14pt}{6pt}
\titlespacing*{\subsection}{0pt}{10pt}{4pt}

% --- Encabezado / pie -------------------------------------------------------
\pagestyle{fancy}
\fancyhf{}
\renewcommand{\headrulewidth}{0.4pt}
\fancyhead[L]{\footnotesize\color{muted} Guion -- Sueño y memoria}
\fancyhead[R]{\footnotesize\color{muted} TP Neurociencias del Desarrollo Productivo}
\fancyfoot[C]{\footnotesize\color{muted}\thepage}

% --- Cajas (barra lateral + tinte via colortbl) -----------------------------
\newcommand{\cajaW}{\dimexpr\linewidth-3pt-2\tabcolsep-1pt\relax}
\newenvironment{narracion}{%
  \par\smallskip\noindent
  \begin{tabular}{@{}!{\color{navy}\vrule width 3pt}>{\columncolor{navy!6}\color{ink}}p{\cajaW}@{}}}%
  {\end{tabular}\par\smallskip}
\newenvironment{clavebox}{%
  \par\smallskip\noindent
  \begin{tabular}{@{}!{\color{amber}\vrule width 3pt}>{\columncolor{amber!10}\color{ink}}p{\cajaW}@{}}}%
  {\end{tabular}\par\smallskip}
\newenvironment{avisobox}{%
  \par\smallskip\noindent
  \begin{tabular}{@{}!{\color{warn}\vrule width 3pt}>{\columncolor{warn!8}\color{ink}}p{\cajaW}@{}}}%
  {\end{tabular}\par\smallskip}

\newcommand{\pregunta}[1]{\vspace{4pt}\noindent\textbf{\color{navy}P. #1}\par\nopagebreak}
\newcommand{\respuesta}[1]{\noindent\textbf{R.} #1\par}
\newcommand{\slidehdr}[2]{\subsection*{\color{steel}Slide #1 \,\textbar\, #2}\addcontentsline{toc}{subsection}{Slide #1 -- #2}}

\setlist[itemize]{leftmargin=1.4em, topsep=2pt, itemsep=1.5pt}
\setlist[enumerate]{leftmargin=1.6em, topsep=2pt, itemsep=2pt}

% =============================================================================
\begin{document}

% ---------- PORTADA ----------------------------------------------------------
\thispagestyle{empty}
\begin{center}
  \vspace*{2.0cm}
  {\color{amber}\rule{4cm}{2pt}}\\[10pt]
  {\Huge\bfseries\color{navy} Guion de presentación}\\[10pt]
  {\LARGE\color{ink} Sueño, memoria y análisis de señales fisiológicas}\\[8pt]
  {\large\color{muted} Consolidación de memoria declarativa durante el sueño}\\[24pt]
  {\large Trabajo práctico -- Neurociencias del Desarrollo Productivo}\\[10pt]
  {\large\textbf{Integrantes:} Keoni Dubovitsky \,\textbullet\, Joaquín Pelufo \,\textbullet\, Sergio Tenoch Sil Cruz}\\[6pt]
  {\large\color{muted} Fecha de exposición: 18/6/2026}\\[20pt]
  \colorbox{navy!6}{\parbox{0.80\textwidth}{\medskip\centering\normalsize
    \textbf{Tesis oral.} El sueño mejora la consolidación de la memoria declarativa.
    Con los datos del práctico la hipótesis queda \textbf{apoyada de forma descriptiva},
    no demostrada estadística ni causalmente.\medskip}}
  \vfill
  {\small\color{muted} Documento de exposición -- versión 3.\\
   Acompaña a \texttt{2026-06-11\_presentacion-sueno-memoria\_ultrabeamer\_v3.pdf}.
   Duración objetivo: 12--15 minutos.}
\end{center}
\newpage

\tableofcontents
\newpage

% ---------- 1. OBJETIVO ------------------------------------------------------
\section{Objetivo y alcance}
El objetivo del trabajo es estudiar la relación entre el \textbf{sueño} y la
\textbf{consolidación de la memoria declarativa}, integrando dos tipos de datos: una
\textbf{tarea de palabras} (memoria) y un \textbf{registro fisiológico} de sueño
(EEG/EOG/EMG). El práctico se usa como punto de partida para pensar críticamente la
relación entre sueño, memoria, actividad fisiológica y las variables que modulan el
descanso y la consolidación.

\begin{avisobox}
\textbf{Restricción metodológica central.} El sujeto propio registrado con OpenBCI
\textbf{no concilió el sueño}. Por indicación de la cátedra, el análisis fisiológico de
sueño se hace sobre un \textbf{dataset secundario} de una persona que sí durmió. Por eso
\textbf{no se mezclan} ambas fuentes como si fueran el mismo sujeto, ni se correlaciona
directamente el EEG con la tarea de memoria.
\end{avisobox}

\subsection{Diseño del experimento}
El experimento tiene tres etapas y \textbf{dos condiciones}:
\begin{enumerate}
  \item \textbf{TR} -- etapa inicial: todos los participantes aprenden y recuerdan las
        \textbf{asociaciones de palabras} (línea de base).
  \item \textbf{Intervalo} -- acá se separan las condiciones: la condición \textbf{sueño}
        realiza una \textbf{siesta prevista de $\sim$90 min} (etapa de consolidación), mientras que el
        \textbf{grupo control} permanece en \textbf{vigilia} (despierto).
  \item \textbf{TS} -- test final: ambos vuelven a evaluar el recuerdo, y se compara
        el desempeño final y el cambio TR$\to$TS entre condiciones.
\end{enumerate}
El contraste \textbf{sueño vs.\ control en vigilia} (y el cambio TR$\to$TS) es lo que
en un diseño controlado permitiría estimar el efecto del sueño. En este TP se interpreta
de forma \textbf{descriptiva}, porque la condición sueño tiene $n=1$.

% ---------- 2. HIPOTESIS -----------------------------------------------------
\section{Hipótesis y predicciones}
\textbf{Hipótesis principal.} El sueño mejora la consolidación de la memoria declarativa.

\begin{itemize}
  \item \textbf{Predicción conductual:} la condición con sueño debería mostrar mayor
        retención TR$\to$TS o mejor desempeño final posterior que el grupo control en vigilia.
  \item \textbf{Predicción fisiológica:} el sueño debería mostrar arquitectura NREM,
        idealmente con sueño de ondas lentas (SWS), y ritmos relevantes para la consolidación.
  \item \textbf{Predicción mecanística:} ondas lentas, husos de sueño y ripples
        hipocampales sostienen la consolidación activa.
\end{itemize}

\begin{clavebox}
\textbf{Criterio de decisión (definido a priori).} La hipótesis se considera
\emph{apoyada} si el patrón conductual y la arquitectura fisiológica son compatibles con
ella. No se busca ``demostrar'': el tamaño muestral es pequeño ($n=1$ en condición sueño)
y las fuentes de memoria y EEG están separadas. El lenguaje correcto es ``compatible con''
o ``apoya descriptivamente'', nunca ``queda demostrado''.
\end{clavebox}

% ---------- 3. QUE ESPERABAMOS ENCONTRAR -------------------------------------
\section{Qué esperábamos encontrar y con qué criterio analizamos}
Esta sección explicita, antes de mostrar resultados, \textbf{qué anticipábamos observar}
y \textbf{qué decisiones metodológicas} tomamos y por qué. Sirve para mostrar que el
análisis siguió un criterio y no se ajustó a los resultados.

\subsection{Lo que esperábamos encontrar}
\begin{itemize}
  \item \textbf{En la tarea (conductual):} mayor retención TR$\to$TS o mejor desempeño final
        en la condición sueño que en el control en vigilia, sobre todo en la recuperación más
        exigente (la forma exacta de las palabras).
  \item \textbf{En el hipnograma:} un descenso ordenado de vigilia a NREM. Como el registro
        estaba planificado como siesta de alrededor de 90 minutos, esperábamos \emph{poder} llegar al
        sueño profundo (SWS) e incluso a un primer episodio de \textbf{REM}.
  \item \textbf{En el espectro:} potencia \textbf{delta} máxima en SWS, presencia de
        \textbf{potencia sigma compatible con husos} (12--15 Hz) en S2 y más \textbf{alfa/beta} en vigilia.
  \item \textbf{En el mecanismo:} ritmos compatibles con la consolidación activa.
\end{itemize}

\begin{clavebox}
\textbf{Qué encontramos realmente.} El registro \textbf{sí} llegó a S3/SWS (bloque
sostenido de 45,5 a 68 min), pero \textbf{no mostró S4 ni REM}. Por eso el análisis se
centra en la comparación entre fases NREM (Vigilia, S1, S2, S3) y en el sueño profundo,
sin forzar una comparación con REM que los datos no permiten.
\end{clavebox}

\subsection{Qué criterio utilizamos y por qué}
\begin{enumerate}
  \item \textbf{Separación de fuentes.} La memoria (planilla, 4 sujetos) y el EEG (dataset
        secundario) se analizan por separado. \emph{Por qué:} el sujeto propio no durmió y
        la cátedra indicó usar datos secundarios; mezclarlos induciría conclusiones
        causales inválidas.
  \item \textbf{Interpretación del scoring.} El archivo solo contiene los códigos
        \textbf{0, 1, 2 y 3} = Vigilia, S1, S2 y \textbf{S3/SWS}, un descenso NREM
        consecutivo. El registro \textbf{llegó hasta S3} (sueño profundo); no hay S4 ni
        REM. \emph{Por qué:} es una lectura fisiológicamente coherente con la nota de cátedra
        (``solo llegó a fase 3'') y con la fisiología; usamos solo la primera columna del
        \texttt{.txt} y la segunda se ignora. Se presenta como scoring descriptivo, no como
        estadificación clínica definitiva.
  \item \textbf{Filtrado antes de medir.} Pasa banda 0,3--35 Hz + notch 50 Hz. \emph{Por
        qué:} la deriva/DC y el ruido de línea distorsionan el espectro; sin filtrar, la
        potencia delta aparecía inflada en todas las fases.
  \item \textbf{Épocas de 30 s sin solapamiento.} \emph{Por qué:} el scoring fue hecho en
        épocas de 30 s sin solapamiento; las ventanas de análisis se alinean 1:1.
  \item \textbf{Control de calidad de canal.} Comparamos C3 y C4. \textbf{Descartamos C4}
        (desvío de 513 uV, $\sim$18 veces el de C3, con saturación) y analizamos
        \textbf{C3}. \emph{Por qué:} promediar C4 contaminaba el espectro e invertía el
        patrón de delta entre fases.
  \item \textbf{Sigma como proxy de husos.} Medimos la banda sigma 12--15 Hz, pero
        \emph{no} corrimos un detector formal de husos. \emph{Por qué:} sin un método de
        detección no afirmamos ``hay husos''; hablamos de la banda sigma como proxy y
        proponemos el detector como trabajo futuro.
  \item \textbf{Lenguaje prudente.} Reportamos resultados descriptivos y separamos
        inferencias causales. \emph{Por qué:} $n=1$ en sueño y fuentes separadas no
        permiten estadística ni causalidad.
\end{enumerate}

% ---------- 4. RESUMEN POR SLIDE ---------------------------------------------
\section{Resumen ejecutivo por slide}
\renewcommand{\arraystretch}{1.25}
\begin{longtable}{@{}p{0.7cm}p{6.1cm}p{8.5cm}@{}}
\toprule
\textbf{\#} & \textbf{Slide} & \textbf{Mensaje central (claim)} \\
\midrule
0 & Portada & Pregunta, hipótesis e integrantes. \\
1 & Resumen & La evidencia apoya descriptivamente la hipótesis, no la demuestra. \\
2 & Introducción & La memoria declarativa necesita consolidarse; el sueño es la ventana. \\
3 & Marco teórico aplicado & Consolidación activa, ondas lentas, sigma/husos y ripples: qué buscamos en el TP. \\
4 & Hipótesis & Hipótesis y predicciones contrastables + criterio de decisión. \\
5 & Diseño del experimento & TR $\to$ (siesta prevista / control vigilia) $\to$ TS. \\
6 & Procedencia de datos & Dos fuentes (memoria y EEG) que no se mezclan. \\
7 & Métodos conductuales & Recuerdo formal vs.\ semántico; condición sueño $n=1$. \\
8 & Métodos EEG & Filtrar antes de medir; C4 descartado, se usa C3. \\
9 & Scoring & El registro llegó hasta S3/SWS; sin S4 ni REM. \\
10 & Resultados memoria & Desempeño final alto en sueño y cambio TR$\to$TS descriptivo. \\
11 & Hipnograma & Descenso NREM completo hasta S3/SWS, sin REM. \\
12 & Fases y latencias & Predomina NREM; $\sim$28\% en sueño profundo. \\
13 & Potencia por banda & Delta crece hacia S3; sigma, proxy compatible con husos, pico en S2. \\
14 & Sigma y memoria & Integramos el mecanismo; no inventamos una correlación. \\
15 & Variables externas & Cronotipo, cafeína, estrés, pantallas, ambiente. \\
16 & Discusión & Lo que encaja y las limitaciones, sin sobre-afirmar. \\
17 & Conclusión & El sueño como aliado de la memoria, con cautela. \\
\bottomrule
\end{longtable}

% ---------- 5. GUION DETALLADO -----------------------------------------------
\section{Guion detallado por slide (texto sugerido)}

\slidehdr{0}{Portada}
\begin{narracion}
``Buenos días. Somos Keoni, Joaquín y Sergio. Vamos a presentar nuestro trabajo sobre
sueño, memoria y señales fisiológicas. La pregunta que organiza todo es simple: ¿el sueño
mejora la consolidación de la memoria declarativa?''
\end{narracion}

\slidehdr{1}{Resumen}
\begin{narracion}
``En 45 segundos: evaluamos descriptivamente la relación entre sueño, memoria declarativa y
señales fisiológicas. Aprendimos asociaciones de palabras, comparando una condición que durmió
una siesta con un grupo control que se mantuvo en vigilia. El EEG de sueño viene de un
dataset secundario de una persona que sí durmió. El sujeto en condición sueño tuvo mayor
desempeño final y una mejora TR a TS; el EEG muestra una arquitectura NREM con sueño profundo.
Conclusión cauta: la evidencia apoya la hipótesis de forma descriptiva, pero no la confirma,
porque memoria y EEG son fuentes distintas y el grupo sueño tiene un solo sujeto.''
\end{narracion}

\slidehdr{2}{Introducción}
\begin{narracion}
``La memoria declarativa son hechos y asociaciones que podemos expresar; se aprende rápido
pero queda frágil y necesita consolidarse. El sueño es la ventana donde esa consolidación
es más eficaz. En este TP lo ponemos a prueba: aprendemos asociaciones de palabras y
comparamos el recuerdo con y sin sueño.''
\end{narracion}

\slidehdr{3}{Marco teórico aplicado}
\begin{narracion}
``No vamos a re-explicar toda la teoría, sino a conectarla con el TP. La consolidación
activa dice que en sueño profundo se reactivan las memorias y se transfieren del hipocampo
a la neocorteza, gracias al acople de ondas lentas, husos y ripples. La homeostasis
sináptica agrega que el NREM mejora la relación señal/ruido. Por eso, en el TP buscamos tres
cosas concretas: mejor retención tras dormir, presencia de sueño profundo en el hipnograma, y
delta y sigma compatible con husos en el EEG.''
\end{narracion}

\slidehdr{4}{Hipótesis y predicciones}
\begin{narracion}
``Nuestra hipótesis: el sueño mejora la consolidación declarativa. De ahí tres predicciones:
conductual (mayor retención o mejor desempeño final en la condición sueño que en el control),
fisiológica (arquitectura NREM con SWS) y mecanística (ondas lentas, husos y ripples). Fijamos el
criterio de antemano: la hipótesis se apoya si los patrones son compatibles; no la
'demostramos' por el tamaño muestral y la separación de fuentes.''
\end{narracion}

\slidehdr{5}{Diseño del experimento}
\begin{narracion}
``El diseño tiene tres etapas y dos condiciones. Primero, la etapa TR: todos aprenden y
recuerdan las asociaciones de palabras, como línea de base. Después viene el intervalo, y
acá se separan las condiciones: una persona realiza una siesta prevista de alrededor de
90 minutos, que es la etapa de consolidación, mientras que el grupo control se queda
despierto, en vigilia. Por último, la etapa TS: todos vuelven a evaluar el recuerdo.
En un diseño ideal, comparar sueño contra vigilia permitiría estimar el efecto del sueño;
en nuestro TP lo interpretamos descriptivamente por el tamaño muestral.''
\end{narracion}

\slidehdr{6}{Procedencia de datos}
\begin{narracion}
``Tenemos dos fuentes que no se mezclan. Fuente A: la planilla de la tarea de palabras, con
cuatro sujetos, tres en vigilia como grupo control y uno en sueño. Fuente B: un registro
BrainVision de una persona que sí durmió, que la cátedra nos dio como dataset secundario.
Importante: nuestro sujeto propio con OpenBCI no logró dormir, así que ese caso lo usamos
como limitación y tema de discusión, no como dato de sueño.''
\end{narracion}

\slidehdr{7}{Métodos conductuales}
\begin{narracion}
``En TR y TS cada sujeto trabaja con 20 asociaciones de una palabra rara con su definición.
Medimos dos cosas por separado: la palabra objetivo, que es la coincidencia formal exacta, y
la definición, con un scoring semántico. Separar ambas nos deja ver que recuperar la forma
exacta es más exigente que recuperar el significado. El análisis es descriptivo: la condición
sueño tiene un solo sujeto.''
\end{narracion}

\slidehdr{8}{Métodos EEG}
\begin{narracion}
``El EEG lo procesamos con MNE. Antes de medir, filtramos: un notch de 50 Hz y un pasa banda
de 0,3 a 35 Hz, para sacar deriva y ruido. Cortamos en épocas de 30 segundos alineadas al
scoring. Y revisamos la calidad de canal: C4 tenía artefactos enormes y saturación, así que
lo descartamos y trabajamos con C3. Este control es clave: promediar C4 distorsionaba todo el
espectro e invertía el patrón de delta.''
\end{narracion}

\slidehdr{9}{Scoring de sueño}
\begin{narracion}
``El archivo de scoring tiene 159 épocas de 30 segundos, con códigos del 0 al 3. Para este
análisis asumimos la lectura vigilia, S1, S2 y S3, que es el sueño profundo o de ondas lentas.
Lo presentamos como scoring descriptivo, no como estadificación clínica. El registro útil
analizado dura 79,5 minutos; aunque el protocolo apuntaba a una siesta cercana a 90 minutos,
no aparecen S4 ni REM. Así que comparamos las fases por profundidad, con foco en S3.''
\end{narracion}

\slidehdr{10}{Resultados conductuales}
\begin{narracion}
``El resultado conductual hay que leerlo en dos niveles. En desempeño final, el sujeto en
sueño obtuvo 17 de 20, frente a un promedio de 5 de 20 en vigilia. Pero para consolidación
también miramos el cambio desde la línea de base: sujeto 1 vigilia 0 a 0, sujeto 2 vigilia
10 a 14, sujeto 3 vigilia 1 a 1 y sujeto 4 sueño 15 a 17. Entonces el sujeto sueño tuvo el
TS más alto y una mejora de +2, pero ya partía alto. Por eso lo presentamos como patrón
compatible con la hipótesis, no como prueba causal.''
\end{narracion}

\slidehdr{11}{Hipnograma}
\begin{narracion}
``Este es el hipnograma de la persona que durmió. Se ve un descenso NREM completo: de vigilia
a S1, S2 y un bloque sostenido de S3, el sueño profundo, entre los minutos 45 y 68. La
latencia de sueño fue de 4 minutos y a S3 de 45 minutos. No hay REM.''
\end{narracion}

\slidehdr{12}{Tiempo por fase y latencias}
\begin{narracion}
``En la distribución predomina el NREM: S2 es la fase más larga con casi 41\% del tiempo, y el
sueño profundo S3 ocupa cerca del 28\%. La eficiencia de sueño es alta, 86\% del registro. Es
una arquitectura compatible con una siesta con buena proporción de NREM profundo.''
\end{narracion}

\slidehdr{13}{Potencia por banda y fase}
\begin{narracion}
``Acá está la firma espectral. La potencia delta crece con la profundidad: de 57\% en vigilia
a 92\% en S3. La banda sigma absoluta tiene su pico claro en S2, justo donde la teoría ubica
los husos de sueño. Como no aplicamos un detector formal, lo formulamos como potencia sigma
compatible con husos, no como husos confirmados.''
\end{narracion}

\slidehdr{14}{Sigma y memoria}
\begin{narracion}
``Y acá la honestidad metodológica. Lo que sí podemos decir: este registro muestra potencia
sigma compatible con husos en S2 y ondas lentas en S3, ritmos que la teoría liga a la consolidación. Lo que no podemos
decir: que los husos de este EEG expliquen la memoria de esos sujetos, porque son fuentes
distintas. Por eso integramos el mecanismo a nivel teórico, no calculamos una correlación
directa. Además, la sigma es un proxy: un detector formal de husos queda como trabajo futuro.''
\end{narracion}

\slidehdr{15}{Variables externas}
\begin{narracion}
``¿Por qué nuestro sujeto pudo no dormir? Lo discutimos con la teoría: el cronotipo y el
proceso circadiano (una siesta a las cuatro de la tarde podría ir a contra-fase), el café que
podría bloquear adenosina y reducir sueño profundo, el estrés del recuperatorio previo, las pantallas y el
ambiente no habitual con electrodos incómodos. Algunas las reportó el sujeto; otras son
plausibles. Sin mediciones directas, no son causas confirmadas.''
\end{narracion}

\slidehdr{16}{Discusión integrada}
\begin{narracion}
``Integrando todo: lo que encaja es el mayor desempeño final en sueño con una mejora TR a TS
descriptiva, la arquitectura NREM con sueño profundo y potencia sigma elevada en S2, y la
coherencia con consolidación activa y homeostasis sináptica. Las
limitaciones son el $n=1$, las fuentes separadas, que el registro no llegó a S4 ni REM, y la
sigma como proxy. La idea es integrar el mecanismo de la materia sin sobre-afirmar lo que los
datos no permiten.''
\end{narracion}

\slidehdr{17}{Conclusión}
\begin{narracion}
``Cerramos con una respuesta directa: los resultados son compatibles con la hipótesis. El sujeto
en sueño mostró el mayor desempeño final y una mejora descriptiva, y el EEG secundario mostró
NREM con sueño profundo y sigma elevada en S2. El sueño aparece como un aliado de la memoria
declarativa, pero lo afirmamos con la cautela que los datos exigen: apoyo descriptivo y parcial,
no demostración. Y el sujeto que no durmió, lejos de
arruinar el trabajo, nos dio la mejor discusión sobre cronotipo, cafeína, estrés y ambiente.
Gracias.''
\end{narracion}

% ---------- 6. PREGUNTAS -----------------------------------------------------
\section{Preguntas esperables de la profesora y respuestas}

\pregunta{¿La hipótesis se acepta o se rechaza?}
\respuesta{Se \textbf{apoya de forma descriptiva}. Los datos son compatibles (mayor desempeño
final y mejora TR$\to$TS en sueño + arquitectura NREM con sueño profundo en el EEG secundario),
pero no se puede \emph{confirmar} estadística ni causalmente: la condición sueño tiene $n=1$ y
el EEG proviene de otra fuente.}

\pregunta{¿Había grupo control?}
\respuesta{Sí. El diseño tiene dos condiciones: la condición \textbf{sueño} (siesta prevista de $\sim$90 min)
y un \textbf{grupo control en vigilia} (3 sujetos despiertos). Ambos hacen TR y TS; la
comparación permite una lectura descriptiva del efecto esperado del sueño, no una prueba causal fuerte.}

\pregunta{¿Por qué no pueden correlacionar directamente husos y memoria?}
\respuesta{Porque el EEG (dataset secundario de quien durmió) y la tarea de memoria (otros
sujetos) \textbf{no pertenecen al mismo sujeto ni sesión}. Una correlación directa supondría
que son la misma persona. Integramos el vínculo husos--memoria a nivel \textbf{teórico}.}

\pregunta{¿Qué significa SWS y por qué importa?}
\respuesta{SWS es el sueño de ondas lentas, equivalente a S3/N3, la fase más profunda del NREM,
dominada por actividad delta de alta amplitud. Importa porque es donde la teoría de
consolidación activa ubica la \textbf{reactivación hipocampo-cortical} de las memorias
declarativas. En nuestro registro aparece como un bloque sostenido de 45,5 a 68 minutos.}

\pregunta{¿Qué aporta la banda sigma?}
\respuesta{La banda sigma (12--15 Hz) es el rango de los \textbf{husos de sueño}, típicos de S2,
asociados a plasticidad y consolidación declarativa. En nuestros datos la sigma absoluta tiene
su \textbf{pico en S2}, lo esperable. La usamos como \textbf{proxy descriptivo}: no aplicamos un
detector formal de husos.}

\pregunta{¿Están comparando desempeño final o consolidación TR$\to$TS?}
\respuesta{Mostramos ambos. El desempeño final favorece al sujeto sueño (17/20), pero el cambio
TR$\to$TS fue +2 porque ya partía de 15/20. Por eso hablamos de desempeño final alto y mejora
descriptiva compatible con la hipótesis, no de una demostración causal de mayor consolidación.}

\pregunta{¿Por qué el sujeto propio no durmió?}
\respuesta{No lo sabemos con certeza porque no lo medimos. Lo discutimos como convergencia
plausible de factores: fase circadiana y cronotipo (siesta a las 16:00), cafeína, estrés del
recuperatorio previo, pantallas y un ambiente no habitual con electrodos incómodos. Son
\textbf{variables moduladoras plausibles}, no causas confirmadas.}

\pregunta{¿Qué rol tienen cronotipo, estrés y cafeína?}
\respuesta{Son moduladores del dormir y de la consolidación. El \textbf{cronotipo} y el proceso
circadiano determinan cuándo es fácil dormir. La \textbf{cafeína} puede antagonizar la adenosina,
aumentar la latencia y reducir el SWS. El \textbf{estrés} puede elevar activación fisiológica,
dificultar conciliar y perjudicar la consolidación declarativa en el sueño temprano.}

\pregunta{¿Qué limitaciones tiene el scoring?}
\respuesta{Trabajamos solo con la primera columna del archivo (la segunda se ignora). El
registro no incluye épocas de S4 ni REM, así que la comparación NREM vs REM no es posible.
Hablamos de fases según una lectura asumida del scoring provisto, no de una polisomnografía
clínica formal. Si la cátedra pide la clave exacta, la respuesta correcta es que la codificación
debe confirmarse oficialmente.}

\pregunta{¿Por qué usar C3/C4?}
\respuesta{Son derivaciones \textbf{centrales}, estándar para la estadificación del sueño: captan
bien husos y ondas lentas. En nuestro caso \textbf{C4 estaba artefactado} (desvío $\sim$18 veces
mayor y saturación), así que lo descartamos y analizamos \textbf{C3}, con señal limpia.}

\pregunta{¿Qué cambiarían en un próximo experimento?}
\respuesta{(1) Registrar EEG y memoria en el \textbf{mismo sujeto} para poder correlacionar.
(2) Asegurar condiciones de sueño (horario, ambiente, control de cafeína). (3) Aumentar el
\textbf{n}. (4) Aplicar un \textbf{detector formal de husos} y de ondas lentas. (5) Documentar
cronotipo y variables contextuales.}

% ---------- 7. CIERRE METODOLOGICO -------------------------------------------
\section{Cierre metodológico (recordatorio para nosotros)}
\noindent{\bfseries\color{warn}Reglas que no se negocian durante la exposición:}
\begin{itemize}[leftmargin=1.5em]
  \item No afirmar que el sujeto propio OpenBCI durmió.
  \item No presentar el dataset secundario como si fuera el mismo sujeto o sesión.
  \item Recordar que la memoria en condición sueño tiene $n=1$.
  \item No calcular correlación directa husos--memoria con fuentes separadas.
  \item La sigma se integra como proxy compatible con husos y análisis descriptivo.
  \item Cafeína, estrés, pantallas, cronotipo y ambiente son variables plausibles, no causas
        demostradas.
  \item Mantener el lenguaje: ``compatible con'', ``apoya descriptivamente''.
\end{itemize}

% ---------- 8. FUENTES -------------------------------------------------------
\section{Material de la cátedra usado}
\begin{itemize}
  \item Clase 3 -- Sueño y memoria (consolidación activa, oscilaciones, homeostasis sináptica,
        polisomnografía y bandas EEG).
  \item Clase 6 -- Estrés (eje HPA, cortisol, efectos sobre memoria y sueño).
  \item Clase 7 -- Ritmos circadianos e higiene del sueño (Proceso S y C, cronotipo).
  \item Clase 8 -- Fármacos (cafeína y adenosina; acetilcolina y consolidación).
  \item Clase 9 -- EEG (polisomnografía: EEG/EOG/EMG, ritmos, fases, hipnograma).
  \item Enunciado oficial del trabajo práctico.
\end{itemize}

\end{document}
